\section{评测是终点，也是起点}

访谈最后强调，机器人和大模型都面临评测卡口。没有好评测，就无法知道模型是否真的变聪明；没有可规模化评测，就无法形成数据闭环。评测不是训练结束后的报告，而是决定下一轮数据采集和训练方向的起点。

\subsection{评测的三种角色}

\noindent\begin{longtable}{p{0.22\textwidth}p{0.34\textwidth}p{0.35\textwidth}}
\toprule
角色 & 功能 & 例子 \\
\midrule
裁判 & 判断当前模型是否更强 & benchmark、端到端成功率、真实任务通过率。 \\
数据挖掘器 & 暴露失败和长尾场景 & 失败 case 自动回流训练集。 \\
方向盘 & 决定下一步训练和采集方向 & 哪类任务短板最大，就补哪类数据。 \\
\bottomrule
\end{longtable}

\begin{importantbox}{评测也是数据资产}
高质量评测定义了“什么叫进步”。当模型越来越强，评测本身会成为稀缺数据，因为它需要更懂任务、更懂失败、更懂未来需求的人来设计。
\end{importantbox}

\subsection{本章小结}

数据闭环必须以评测为中心。没有评测，就没有方向；没有方向，数据只会堆积而不会变成能力。

